\documentclass[11pt]{article}
\usepackage[margin=0.9in]{geometry}
\usepackage{amsmath,amssymb,amsthm,booktabs,mathrsfs}
\usepackage[T1]{fontenc}
\usepackage{lmodern,microtype}
\usepackage[colorlinks=true,linkcolor=blue,urlcolor=blue]{hyperref}
\newtheorem{theorem}{Theorem}[section]
\newtheorem{lemma}[theorem]{Lemma}
\newtheorem{proposition}[theorem]{Proposition}
\newtheorem{corollary}[theorem]{Corollary}
\theoremstyle{remark}
\newtheorem{remark}[theorem]{Remark}
\newcommand{\norm}[1]{\left\|#1\right\|}
\title{Toward a four-dimensional Yang--Mills construction\\
\large Notebook 24: inherited interactions and the next-step norm test}
\author{Research working document}
\date{17 October 2026}
\begin{document}
\maketitle
\begin{abstract}
We separate an inherited analytic quantum correction into its
stationary-background value and its fluctuation increment in
the actual ordinary pivot chart. Directional Cauchy bounds the
increment by a factor of order $gR$ divided by the available
analytic margin, without a carrier-size derivative loss.
The exact pivot incidence confines its projected support to
the original blocks and their relevant endpoint neighbors.
Connected-carrier compression then yields an explicit
sufficient inequality for the stronger interaction norm
required by the next Gaussian integration. Under that
inequality and the existing covariance hypotheses, the
inherited contribution to the connected logarithm has the
same small factor. The inequality is not yet verified for
the actual output decay constants as the block side varies.
This result makes that parameter test, and the independent
domain-stability problem, explicit.
\end{abstract}

Notebooks 21--23 provide the classical-reference integral,
stationary-background transport, and local gauge covariance.
The following argument treats the additional interactions
produced by earlier integrations. It keeps their full
background term outside the integral and controls the
remaining fluctuation effect under a displayed norm budget.

